\chapter{Estructura matemática de la doble proyección}
\label{chap:sintesis-hmt}

La construcción parte del espacio de todos los caminos admisibles (APP), sin
introducir la recta real como soporte inicial. Sus estados elementales son
TRIT orientados y la célula nonádica se extiende mediante coordenadas de
residuo y cociente euclídeo. El cociclo mixto distingue el orden de los
observables aditivo y multiplicativo, mientras que el producto nativo
extiende la multiplicación local a palabras globales.

El sistema de transporte de fase (TPK) determina la dinámica finita asociada
a la célula. Las 104\,976 semillas se factorizan mediante dos firmas
independientes del resultado decimal: 18 firmas aditivas y 26 multiplicativas.
La combinación inyectiva produce 468 estados, cuya proyección ternaria posee
243 valores. La geometría interna del catálogo contiene una microfibra de
cinco regiones sobre \(010211\), con 1008 rutas semilla y cinco paseos cerrados mínimos
de longitud ocho.

Las palabras ternarias asociadas a \(\pi,e,\varphi\) se prolongan a tríadas
decimales mediante aplicaciones enteras. El calendario, las extensiones y el
acarreo constituyen una estructura dodecafásica. Fijado el estado firmado
enriquecido, el vector \(K\), los operadores \(D_3,D_4,Q\), la órbita de
Hadamard y la evaluación arquimediana proporcionan representaciones
reversibles. Su composición reconstruye internamente \(\alpha\) y reproduce,
como reconocimiento metrológico externo, la cadena central publicada en 2018.

En la instancia finita calibrada, el estado dodecafásico se compara con una
proyección excepcional. El vector
\(K\) y los soportes del código recuperan por dos rutas la bandera
\[
B_K\subset H^-_\alpha.
\]
Desde la matriz Paley--Witt se construyen el código ternario extendido y el
diseño \(W_{12}\); las 495 estrellas generan \(M_{12}\). En la rama
reticular, la realización discriminante de \(A_2^{12}\), el origen recuperado
y la cámara radial seleccionan un vector de norma 54. El certificado prueba
que su 3-vecino es par, unimodular y sin raíces; la clasificación clásica lo
identifica con Leech, y los teoremas de Frenkel--Lepowsky--Meurman prolongan
la construcción hacia \(V^\natural\) y el grupo Monstruo. En la rama binaria,
una dodecada del código de Golay y su alineamiento con las doce posiciones
HMT determinan el diseño residual \(W_{12}\) y la elevación única de cada
hexada a una octada de \(W_{24}\). Las dos ramas comparten el código como nodo
inicial y conservan codominios distintos.

La definición de número resultante es relacional. Una historia compatible
determina intervalos cilíndricos; si están encajados y su diámetro tiende a
cero, su intersección contiene un único real. El valor real identifica las
historias de una misma fibra. La representación completa del estado conserva
además la orientación, el acarreo, la coordenada de observable y la clase de
trayectoria. La estructura de cuerpo del cociente se obtiene cuando la suma y
el producto descienden por las fibras. Para el sector canónico de palabras,
el producto nativo de \(\mathcal A_9^\#\) establece el entrelazador
multiplicativo requerido.

Los principales teoremas finitos están acompañados por verificadores
autocontenidos. Estos reconstruyen el catálogo, la imagen de 42 ternas, la
factorización \(468=26\times18\), las cinco regiones y sus ciclos, la
aritmética celular, la reconstrucción de \(\alpha\), el núcleo Witt y el
producto nativo. También certifican el mínimo radial, las condiciones I1--I3
y la ausencia de raíces del 3-vecino. El proyecto Lean verifica de forma
independiente identidades algebraicas seleccionadas.

Las hipótesis determinan asimismo el alcance exacto de los teoremas. El estado
dodecafásico firmado es una premisa enriquecida de la reconstrucción relativa
de \(\alpha\); la cobertura de toda \(\mathbb R\) requiere construir una
historia admisible para cada real; no se emplea una acción directa del
Monstruo sobre palabras decimales; y la naturalidad global de todas las
evaluaciones requiere entrelazadores adicionales. Estas condiciones delimitan
con precisión la extensión global de los resultados demostrados.

El resultado finito identifica una bandera común obtenida por una evaluación
arquimediana y una prolongación excepcional relativa a su calibración. La
extensión global se expresa como el problema de construir una ambientación común,
estable a toda profundidad, y los entrelazadores naturales que prolonguen esa
compatibilidad. Los objetos, aplicaciones, fibras, cociclos, límites inversos
y certificados desarrollados en los capítulos precedentes fijan con precisión
el dominio actualmente construido de esa extensión.
